% !PRINTED 229
\weberpage{229}
If, in addition, one also forms $\varphi\psi$, $\varphi^{-1}\psi$, $\varphi\chi\psi$, $\chi\varphi^{-1}\psi$, then it follows that:
\[
\begin{aligned}
\varphi\psi &= \Theta\chi\psi\Theta^{-2}, \quad \varphi^{-1}\psi = \Theta^{-2}\chi\psi\Theta, \quad \varphi\chi\psi = \Theta^{-1}\chi\psi\Theta^2, \\
&\chi\varphi^{-1}\psi = \Theta^2\chi\psi\Theta^{-1},
\end{aligned}
\]
from which one sees that one would come to no other result if one wished to assume the substitution of the third order from which one starts in the second form $\Theta^r\chi\psi\Theta^s$. The whole group $Q$ is thus completely determined by the assumed four-group $D_2$, and one obtains it in the form
\weq{3}{
\begin{array}{ll}
1 & \psi \\
\varphi=\Theta\chi\Theta^2, & \varphi\psi=\Theta\chi\psi\Theta^{-2}, \\
\varphi^{-1}=\Theta^{-2}\chi\Theta^{-1}, & \varphi^{-1}\psi=\Theta^{-2}\chi\psi\Theta, \\
\chi & \chi\psi \\
\varphi\chi=\Theta^{-1}\chi\Theta^{-2}, & \varphi\chi\psi=\Theta^{-1}\chi\psi\Theta^2, \\
\varphi^{-1}\chi=\Theta^2\chi\psi\Theta^{-1}, & \varphi^{-1}\chi\psi=\Theta^2\chi\Theta.
\end{array}
}

One can derive this group from the substitutions $\varphi, \psi, \chi$ as the generators and put it into the form:
\weq{4}{Q = \varphi^r\chi^s\psi^t, \qquad \begin{gathered}r=0,1,2,\\ s=0,1; \quad t=0,1.\end{gathered}}

That a tetrahedral group is really represented in this way follows from the compositions, which one easily finds from (1) by means of \S.\,58, (23), (25):
\weq{5}{
\begin{aligned}
\psi\varphi &= \varphi\chi\psi, \quad \chi\varphi = \varphi\psi, \\
\psi\varphi^2 &= \chi\varphi^2\psi, \quad \chi\varphi^2 = \varphi^2\chi\psi. \quad \psi\chi = \chi\psi. % [sic: Weber prints a full stop between the two relations, unlike the commas separating the parallel relations; the printed punctuation is preserved.]
\end{aligned}
}

Since five four-groups $D_2$ are contained in the tetrahedral group, the icosahedral group has five tetrahedral groups as divisors. % [sic: Weber prints Tetraëdergruppe; mathematically Ikosaëdergruppe is intended.]
These we can derive from $Q$ by transformation by means of the powers of $\Theta$, and we obtain them in the form
\[
\Theta^{-r}Q\Theta^r \quad r=0,1,2,3,4.
\]

Apart from the identity, these groups have no substitution common to all of them. For the two groups $Q$ and $\Theta^{-1}Q\Theta$ have in common with one another, apart from the identity, only the two elements
\[
\Theta^{-2}\chi\Theta^{-1}, \quad \Theta\chi\Theta^2
\]
and of these neither occurs in $\Theta Q\Theta^{-1}$.

From this it now follows, by Theorem 2. in \S.\,28, that the icosahedral group is isomorphic to a permutation group
%CONT>>
